% concatenated from Gallardo-Partington_August_2026.tex
\documentclass[11pt]{amsart}
\usepackage{amssymb}
\usepackage{amscd}
\usepackage{color}
\usepackage{comment}
\usepackage[dvipsnames]{xcolor}
\usepackage{url}
\usepackage{hyperref}
\usepackage{enumitem}





%%%%%%%%%%%%%%%%%%%%%
\addtolength{\textheight}{.4in}
\addtolength{\topmargin}{-.2in}
\addtolength{\textwidth}{1.6in}
\addtolength{\oddsidemargin}{-.7in}
\addtolength{\evensidemargin}{-.7in}
%\addtolength{\itemsep}{1cm}
\parindent=12pt
%%%%%%%%%%%%%%%%%%%%%

\hyphenation{o-pe-ra-tor}
\hyphenation{a-na-ly-sis}

\renewcommand{\baselinestretch}{1,2}

% \usepackage{epstopdf}
% \usepackage{MnSymbol}
\allowdisplaybreaks[4]
%\pagestyle{myheadings}

 \newtheorem{thm}{Theorem}[section]
 \newtheorem*{theorem*}{Theorem}
 \newtheorem{cor}[thm]{Corollary}
 \newtheorem{lem}[thm]{Lemma}
 \newtheorem{lemma}[thm]{Lemma}
 \newtheorem{prop}[thm]{Proposition}
 \newtheorem{proposition}[thm]{Proposition}
 \theoremstyle{definition}
 \newtheorem{defn}[thm]{Definition}
\theoremstyle{remark}
 \newtheorem{rem}[thm]{Remark}
\newtheorem{remark}{Remark}
\newtheorem{exm}[thm]{Example}
 \numberwithin{equation}{section}
%\def{\beginpf}{\begin{proof}}
%\def{\endpf}{\end{proof}}
\newcommand{\la}{\langle}
\newcommand{\ra}{\rangle}
\newcommand{\ol}{\overline}
\newcommand{\ma}{\mathcal}
\newcommand{\Ker}{\mathop{\rm ker}\nolimits}
\newcommand{\re}{\mathop{\rm Re}\nolimits}



\newenvironment{demode}
{\noindent {{\it Proof of }}}%
{\par \hfill }





\def\CC{{\mathbb C}}
\def\C{\CC} % Bill's choice
\def\DD{{\mathbb D}}
\def\D{\DD}
\def\dim{\mathop{\rm dim}\nolimits}
\def\lcm{\mathop{\rm lcm}\nolimits}
%\def\Id{{\rm Id}}
\def\Id{I} % we seem to use both Id and I
\def\LL{\mathcal L}
\def\EE{\mathcal E}
\def\NN{\mathbb N}
\def\RR{\mathbb R}
\def\T{\mathcal T}
\def\TT{\mathbb T}
\def\beq{\begin{equation}}
\def\eeq{\end{equation}}
\def\Res{\mathop{\rm Res}\nolimits}
\def\ds{\displaystyle}
\def\diag{\mathop{\rm diag}\nolimits}
\def\op{\|_{\rm op}}
\def\spam{\mathop{\rm span}\nolimits}
\def\rang{\mathop{\rm ran}\nolimits}
\def\phi{\varphi}
\newcommand{\x}{\textnormal{\textbf{x}} }
\newcommand{\uu}{\textnormal{\textbf{u}} }


\newcommand{\clos}{\overline}
\newcommand{\norm}[1]{\left\|#1\right\|}
\newcommand{\ip}[2]{\left\langle #1,#2\right\rangle}
\newcommand{\spn}{\operatorname{span}}



\title{Frame constructions associated with operator orbits}

%   Information for first author
\author{Eva A. Gallardo-Guti\'{e}rrez}
%    Address of record for the research reported here
\address{Eva A. Gallardo-Guti\'errez \newline
Departamento de An\'alisis Matem\'atico y Matem\'atica Aplicada,\newline
Facultad de Matem\'aticas,
\newline Universidad Complutense de
Madrid, \newline
Plaza de Ciencias 3, 28040 Madrid,  Spain
\newline
and Instituto de Ciencias Matem\'aticas ICMAT,
\newline Madrid,  Spain }
\email{eva.gallardo@mat.ucm.es}


%   Information for second author
\author{Jonathan R. Partington}
%    Address of record for the research reported here
\address{Jonathan R. Partington, \newline
School of Mathematics, \newline
University of Leeds, \newline
Leeds LS2 9JT, United Kingdom}
\email{J.R.Partington@leeds.ac.uk}



\thanks{Both authors are partially supported by Plan Nacional  I+D grant no. PID2022-137294NB-I00, \& PID2025-169474NB-C21, Spain. The first author is also supported by
 the Spanish Ministry of Science and Innovation, through the ``Severo Ochoa Programme for Centres of Excellence in R\&D'' CEX2023-001347.}

%    General info
\subjclass[2020]{46B15,30H10}

\date{August 2026}

\keywords{frame, operator orbit, Carleson interpolation, weighted composition operator, Beurling--Lax theory, M\"untz--Szász  condition}



\begin{document}

\begin{abstract}
This paper studies frames in Hilbert spaces generated by the orbits of (in)-finitely many vectors under a single operator, presenting new results on multiplication operators and operators composed of Jordan blocks, which generalizes existing results of Cabrelli, Molter, Paternostro and Philipp  by means of techniques which deal with weighted interpolation, weighted composition operators, and Beurling--Lax theory related to shifts of infinity multiplicity.
Likewise, we discuss Carleson frames and give counterexamples to a recent conjecture of Aldroubi, Cabrelli, Krishtal and Molter.


\end{abstract}

\maketitle


\section{Introduction}
A sequence $(x_n)_{n \ge 0}$  in a infinite dimensional, separable, complex Hilbert space $H$ is said to be a {\em frame\/} if
there are constants $C_1,C_2 > 0$ such that
\[
C_1 \|x\|^2 \le \sum_{n=0}^\infty |\langle x,x_n\rangle|^2 \le C_2 \|x\|^2
\]
for all $x \in H$.
Frames extend the notion of an orthonormal basis: while a basis provides a unique expansion, a frame allows multiple representations of the same vector. This redundancy is often advantageous, as it leads to increased numerical stability and robustness against noise, losses, or perturbations of the frame elements (see, for example, \cite{meyer}).

\smallskip

Associated with any frame there are key operators, such as the analysis operator, synthesis operator, and frame operator, which play a central role in reconstruction. In particular, the frame operator is
defined as the mapping
\[
S:x \mapsto \sum_{n=0}^\infty \langle x,x_n \rangle x_n
\]
and it is
bounded, positive, and invertible, guaranteeing that every vector in $H$ can be recovered from its frame coefficients. Frames include important special cases such as tight frames and Parseval frames, which simplify reconstruction formulas. Recall that the sequence $(x_n)_{n \ge 0}$  is said to be a {\em tight frame\/} if we can take $C_1=C_2$ and a {\em Parseval frame\/} if, in addition, both $C_1$ and $C_2$ are equal to 1; or equivalently,
the \emph{analysis operator}
\[
\begin{array}{ccl}
C_{(x_n)}: H & \longrightarrow & \ell^2 \\[4pt]
x & \longmapsto & \big(\langle x, x_n \rangle\big)_{n\geq 0}
\end{array}
\]
is an isometry. It is well-known (see \cite[Chap.~4]{meyer}) that an equivalent condition for a sequence
$(x_n)_{n \ge 0}$ to be a frame on a Hilbert space $H$  is that the \emph{synthesis operator}
\begin{eqnarray}\label{eq:bddsurj}
C_{(x_n)}^*: \ell^2 & \longrightarrow  & H, \nonumber  \\ [6pt]
(a_n)_{n\geq 0} & \longmapsto & \sum_{n=0}^\infty a_n x_n
\end{eqnarray}
should be bounded and surjective. Note that $S=C_{(x_n)}^*C_{(x_n)}$.

\smallskip

As an important class of frames, \textbf{frames of iterations}, also known as \emph{dynamical frames}, arise from the action of a bounded linear operator $T$ on the Hilbert $H$. The key question here is finding conditions on $T$ and a collection of vectors $(u_j)_{j\in J}\subset H$ such that the family
\[
\{T^n u_j : j \in J,\; n\ge 0\}
\]
forms a frame for $H$. In \cite{CMPP20}, Cabrelli, Molter, Paternostro and Philipp found necessary and sufficient conditions for the existence
of dynamical frames (see also \cite{ACNP25}). In particular, they proved that there exists a collection $(u_j)_{j\in J}\subset H$ such that
$\{T^n u_j : j \in J,\; n \ge 0\}$ is a (dynamical) frame for $H$ if and only if $T$ is similar to a contraction $C$ whose adjoint $C^*$ is strongly stable, namely, there exists a bounded invertible operator $P$ such that
$P^{-1}TP = C$, $\|C\|\le 1$, and $(C^*)^n x \to 0$ for every $x\in H$.
Moreover, in order to generate a Parseval frame of iterations, the authors showed that the equivalence holds when $T$ is indeed a contraction whose adjoint is strongly stable.
Consequently, only contractions can produce Parseval frames of iterations, which turn out to be particularly important since their canonical dual coincides with the frame itself, being relevant in applications. We refer to the recent work \cite{ACKM25} where the authors present a unified survey of dynamical sampling and
its interplay with frame theory as well as the references therein.

\smallskip

Motivated by these results, in this manuscript we study dynamical frames for specific classes of operators. In particular, in Section \ref{sec:2} we are concerned with normal operators with a Riesz basis of eigenvectors, and we relate them to classical and more recent results in interpolation theory, namely, to \emph{weighted interpolation}. At this regard, we remark that the case of diagonal operators acting on $\ell^2$ when a single vector is iterated was explicitly described in \cite{CMPP20}, connecting the problem with interpolation sequences in the Hardy space on the open unit disc $H^2(\DD)$. In particular, for diagonal operators of multiplicity one, the following theorem is proved in \cite{ACMT17}:

\smallskip

\begin{thm}[\cite{ACMT17}]\label{thm:diagonal}
Let $D = \sum_{j\in J} \lambda_j P_j$ be a diagonal operator where $J$ is a finite or countable index set, $(\lambda_j)_{j\in J}\subset \mathbb{C}$ is a
bounded sequence of scalars, and $(P_j)_{j\in J}$ is a sequence of orthogonal
projections satisfying $P_j P_k = 0$ for $j \neq k$ and $\sum_{j\in J} P_j = \mathrm{Id}$. Assume further $\lambda_j \neq \lambda_k$ for $j \neq k$ and
each $P_j$ is rank one and let $\uu =(u_k)_{k\in\mathbb{N}} \in \ell^2$. Then the family
\[
\{D^{\ell} \uu : \ell = 0,1,\ldots\}
\]
is a frame if and only if the following conditions hold:
\begin{enumerate}
\item $|\lambda_k| < 1$ for all $k$;
\item $|\lambda_k| \to 1$;
\item $(\lambda_k)_{k\geq 0}$ satisfies Carleson’s condition
\[
\inf_n \prod_{k \neq n}
\frac{|\lambda_n - \lambda_k|}{|1 - \overline{\lambda_n}\lambda_k|}
\ge \delta,
\]
for some $\delta > 0$;
\item $u_k= m_k \sqrt{1 - |\lambda_k|^2}$ for some sequence $(m_k)_{k\geq 0}$ satisfying
\[
0 < C_1 \le |m_k| \le C_2 < \infty.
\]
\end{enumerate}
\end{thm}

The condition that the eigenvalues have modulus strictly less than one ensures that the operator is contractive and strongly stable. In \cite[Theorem~2.3]{CMPP20} the authors have generalized Theorem~ \ref{thm:diagonal} to arbitrary finite collection $\{\uu_i\}_{i\in I}$ by a rather non-trivial treatment, as they explicitly claimed.


\smallskip

In Section~\ref{sec:3} we deal with \emph{Carleson frames} or \emph{Alcamota frames}, namely frames satisfying the stated conditions (1)-(4).  As pointed out in \cite{ACKM25}, such frames have \emph{excessive redundancy}.
Indeed, the subfamily obtained by selecting each $N$th element from a
Carleson frame remains a frame under a mild additional condition on the spectrum of the diagonal operator $D$ regardless of the choice of $N \in \mathbb{N}$ (see \cite{ChristensenHasannasabPhilippStoeva2024} and \cite{KrishtalMiller2026}). {If $b\in\ell^{2}$, and the orbit $\{D^{\ell}b:\ell=0,1,\dots\}$ forms a Carleson frame for $\ell^{2}$, we first note that not every set $\Lambda=\{\lambda_k\}_{k\geq 0}\subset(0,\infty)$ satisfying the Müntz--Szász divergence condition
\begin{equation}\label{MS condition}
\sum_{\lambda_k\in\Lambda}\frac{1}{\lambda_k}=\infty,
\end{equation}
defines an orbit $\{D^{\lambda}b:\lambda\in\Lambda\}$ that is again a frame for $\ell^{2}$, and we prove that both  \eqref{MS condition} and a separation condition (\emph{a Carleson-type condition}) are necessary to be a frame in $\ell^{2}$. The proof proceeds by passing to a spectral/functional model in which $D$ is represented as multiplication by $t\in(0,1)$ on a Hilbert space $L^2(\nu)$.

\smallskip

In Section~\ref{sec:4} we consider multiplication operators on the Hardy space $H^2(\DD)$ and use weighted composition operators to describe the associated frames.
Finally, Section~\ref{sec:6} is motivated by a result in \cite{CHP20} expressing operator frames in terms of model spaces. We use this to analyse several new classes of frame, and then generalize it to the infinite case.

\section {Normal operators: the weighted interpolation approach} \label{sec:2}

Recall that a \emph{normal operator} on a Hilbert space $H$ is a bounded linear operator $T$ that commutes with its adjoint, that is,
$T T^* = T^* T.$ Equivalently, $T$ is normal if
$\|Tx\| = \|T^*x\| \quad \text{for all } x \in H$. Normal operators include important classes such as self-adjoint and unitary operators,
and they admit a spectral theorem that represents them as multiplication operators on suitable $L^2$-spaces.

\smallskip

Assume $T$ is a normal operator. We start  by pointing out that the authors in \cite{ACCMP17} proved  that there does not exist a Riesz basis of iterations
for a normal operator, that is, for any collection of vectors $G \subset H$, the family of iterates $\{T^n u : u \in G,\; n \ge 0\}$ is no longer a Riesz basis of $H$. Recall that a \emph{Riesz basis} in $H$ is a sequence of vectors
$(f_k)_{k\geq 0}\subset H$ that is the image of an orthonormal basis under a bounded,
invertible linear operator. Equivalently, $(f_k)_{k\geq 0}$ is a Riesz basis if it is complete
in $H$ and there exist constants $0<A\le B<\infty$ such that
\[
A \sum_k |c_k|^2 \;\le\;
\Big\|\sum_k c_k f_k\Big\|^2
\;\le\;
B \sum_k |c_k|^2,
\]
for all finite scalar sequences $(c_k)_{k\geq 0}$. Unlike general frames, a Riesz basis yields unique expansions of every vector in $H$,
while still allowing more flexibility than an orthonormal basis. In particular, Riesz bases are precisely the non-redundant frames.

\medskip

In this regards, the goal is to investigate the existence of redundant frames for normal operators. Obviously, since unitary operators are
never strongly stable, they cannot generate frames of iterations.
\smallskip

Our aim at this regards is making use of weighted interpolation studied by McPhail \cite{mcphail} and generalized in  \cite{JPP07} to provide an extension of \cite[Theorem~2.3]{CMPP20}. For the sake of clarity, we will first deal when the collection $\{\uu_i\}_{i\in I}$ is reduced to just one vector $\uu$.

\smallskip

\subsection*{Interpolation in the scalar case}

Assume $T$ is a diagonalizable normal operator in $\ell^2$  with eigenvalues $(\lambda_k)_{k \ge 0}$ and orthonormal eigenvectors $(\phi_k)_{k \ge 0}$ (in fact, the results we prove will also apply if the eigenvectors form merely a Riesz basis). Let us write $\uu=\sum_{k=0}^\infty u_k \phi_k \in \ell^2$.

\smallskip

Conditions under which $(T^n \uu)_{n \ge 0}$ forms a frame were established in
\cite{ACMT17} and \cite{CMPP20} as mentioned. In order to revisit these results using interpolation
techniques to extend them further, we conveniently work in the Hardy space $H^2(\DD)$. Recall that $H^2(\DD)$ consists of all analytic functions $f$ on the open unit disk $\mathbb{D}$
\[
f(z) = \sum_{n=0}^\infty a_n z^n, \qquad z \in \mathbb{D},
\]
where the sequence of Taylor coefficients $(a_n)_{n\ge 0} \in \ell^2$ and $\|f\|_{H^2}= \|(a_n)_{n\geq 0}\|_{\ell^2}$.


\smallskip

\noindent Hence, mapping the orthonormal basis elements $\phi_n$ to $z^n$, we find that for $f(z)=\sum_{n=0}^\infty a_n z^n\in H^2$,
\[
\sum_{n=0}^\infty  a_n T^n \uu = \sum_{n=0}^\infty a_n \sum_{k=0}^\infty   \lambda_k^n u_k z^k
= \sum_{k=0}^\infty f(\lambda_k)u_k z^k.
\]
Consequently, the boundedness and ontoness of the synthesis operator associated to $(T^n \uu)_{n \ge 0}$, namely, $C_{(T^n \uu)}^*$, is equivalent to the solution of the weighted interpolation problem
\begin{equation}\label{weighted interpolation}
f(\lambda_k)u_k= c_k \qquad (k\geq 0)
\end{equation}
for all sequences $(c_k)_{k\geq 0}$ in $\ell^2$.

\smallskip

The classical interpolation problem asks if there exist a sequence of points $(z_n)_{n\geq 0}$ in $\mathbb{D}$ having the property that, given an arbitrary bounded sequence of complex numbers $(a_n)_{n\geq 0}$, there exists a bounded analytic function $f$ on $\mathbb{D}$ such that
\begin{equation*}
f(z_n) = a_n \qquad (n\geq 0)
\end{equation*}
The existence of such a sequence, called an \emph{interpolating sequence}, was
established independently by L.~Carleson~\cite{Carleson}, W.~Hayman~\cite{Hayman},
and D.~J.~Newman~\cite{Newman}. In fact, Carleson showed that a sequence
$(z_n)_{n\geq 0}$ in $\mathbb{D}$ is an interpolating sequence if and only if there exists
$\delta > 0$ such that
\begin{equation}\label{Carleson}
\prod_{n \neq k}
\left|
\frac{z_k - z_n}{1 - \overline{z}_n z_k}
\right|
\ge \delta
\qquad (k = 0,1,2,\ldots).
\end{equation}
Clearly, such sequences consist of distinct points. For a discussion of independent and related results, see H.~S.~Shapiro and
A.~L.~Shields~\cite{ShapiroShields} and Garnett's monograph \cite{Garnett}.

\smallskip


In 1990, McPhail \cite{mcphail} studied a weighted interpolation problem which, in particular, extended Shapiro and Shield's results. More precisely,  given a sequence $(z_n)_{n\geq 0}$ in $\mathbb{D}$, let $\Lambda = \Lambda_{(z_n)}$ denote the linear map sending an analytic function
$f$ on $\mathbb{D}$ to the sequence of its values on the $z_n$, namely
$(f(z_n))_{n\geq 0}$. In this setting, Carleson’s result says that
$\Lambda (H^\infty) \supset \ell^\infty $, or equivalently, $\Lambda (H^\infty) = \ell^\infty$
if and only if condition~ \eqref{Carleson} holds. Here $H^\infty$ denotes the Hardy space of bounded analytic functions on $\mathbb{D}$ endowed with the sup norm
$\|f\|_{H^\infty}= \sup_{z \in \mathbb{D}} |f(z)|$.


\medskip

Given a sequence $w =(w_n)_{n\geq 0}$ of positive real numbers, let
$\ell^p(w)$ denote the space of sequences $a = (a_n)_{n\geq 0}$ satisfying
\begin{equation}
\|a\|_{\ell^p(w)} =
\begin{cases}
\displaystyle
\left( \sum_{n=0}^\infty |a_n w_n|^p \right)^{1/p} < \infty,
& 1 \le p < \infty, \\
\noalign{\smallskip} \\[8pt]
\displaystyle
\sup_{n \ge 0} |a_n w_n| < \infty,
& p = \infty.
\end{cases}
\end{equation}

The McPhail theorem characterizes the sequences $(z_n)_{n\geq 0}$ in $\mathbb{D}$ such that
$\Lambda_{(z_n)} (H^p(\nu)) = \ell^p(w)$ for some weighted Hardy spaces $H^p(\nu)$ (see \cite{mcphail} for more details).
As a particular instance, the Shapiro and Shields~\cite{ShapiroShields} theorem follows if $\nu = 1$ and $w_n = (1 - |z_n|)^{1/p}$ for  $(n \ge 0)$, that is,
\[
\Lambda_{(z_n)}(H^p) \supset \ell^p(w)
\]
(in fact, equivalently $\Lambda_{(z_n)}(H^p) = \ell^p(w)$) if and only if condition~ \eqref{Carleson} holds.

\smallskip

We now quote a simplified version of the theorem of McPhail \cite{mcphail} for our purposes:


% Let $(w_k)$ be a sequence of positive real numbers.
% We say that a sequence $(d_k)$ lies in  $\ell^2(w)$ if $(d_k w_k)$ lies in $\ell^2$.

\begin{thm}[\cite{mcphail}]\label{thm:mcphail}
Let $w =(w_n)_{n\geq 0}$ be a sequence of positive real numbers and $(z_n)_{n\geq 0}\subset \mathbb{D}$. Then
$\Lambda_{(z_n)}(H^2) = \ell^2(w)$ if and only if   condition~\eqref{Carleson} holds and there exists constants
$C_1, C_2 > 0$ such that
\[
C_1 \le w_n/\sqrt{1-|z_n|} \le C_2
\]
for all $n$.
\end{thm}

With Theorem~\ref{thm:mcphail} at hand, the weighted interpolation problem \eqref{weighted interpolation} is solvable and hence, Theorem \ref{thm:diagonal} follows readily.


\begin{cor}
Assume $T$ is a diagonalizable normal operator in $\ell^2$  with eigenvalues $(\lambda_k)_{k \ge 0}$ associated to a (Riesz) basis of eigenvectors $(\phi_k)_{k \ge 0}$ and let $\uu=\sum_{k=0}^\infty u_k \phi_k \in \ell^2$. The sequence $(T^n \uu)_{n\geq 0}$ forms a frame in $\ell^2$ if and only if the sequence $(\lambda_k)_{k \ge 0}$ is an interpolating sequence in $\DD$ and there exists constants $C_1, C_2 > 0$ such that
\[
C_1 \le |u_k|/\sqrt{1-|\lambda_k|}
\le C_2
\]
for all $k$.
\end{cor}

\subsection*{The multivariable case} Having discussed the scalar case, we deal with the multivariable case. Without loss of generality, we may assume that $T$ is a diagonalizable normal operator in $H^2(\DD)$  with eigenvalues $(\lambda_k)_{k \ge 0}$ associated to a (Riesz) basis of eigenvectors. Let $\{\uu_1,\ldots,\uu_m\}$ be a set in $H^2(\DD)$.  Our aim is studying when
$$\{T^n \uu_k: n \ge 0, k=1,\ldots,m\}$$
is a frame in $H^2(\DD)$.

\smallskip

For such a goal, we will make use of a more general form of McPhail's theorem, given in \cite{JPP07}, pages 537--538 (see also \cite{JZ01}). The aim is solving a vector-valued weighted interpolation problem a little more involved than \eqref{weighted interpolation}, namely,
\begin{equation}\label{weighted interpolation2}
G_k f(\lambda_k)= c_k \qquad (k\geq 0)
\end{equation}
for all sequences $(c_k)_{k\geq 0}$ in $\ell^2$, where the $G_k$ are rank-1 operators.

\smallskip

In order to state the vector-valued weighted interpolation theorem required, one change has to be made, as these works \cite{JZ01,JPP07} take place in the setting of the right-half plane
$$\CC_+=\{ z \in \mathbb{C} : \Re z > 0 \}.$$
Recall that the Hardy space $H^2(\mathbb{C}_+)$ of the right half-plane consists of all functions
$f$ that are analytic in $\mathbb{C}_+$ and satisfy
\[
\|f\|_{H^2(\mathbb{C}_+)}^2=
\sup_{x>0}
\int_{-\infty}^{\infty} |f(x+iy)|^2 \, dy
< \infty.
\]
It is worth remarking that the norms on $H^2(\mathbb{C}_+)$ and $L^2(0,\infty)$ are equivalent via the Laplace
transform. More precisely, the Laplace transform
\[
\mathcal{L}f(s)= \int_0^\infty e^{-st} f(t)\,dt,
\qquad s \in \mathbb{C}_+,
\]
defines (up to a constant) an isometric isomorphism from $L^2(0,\infty)$ onto $H^2(\mathbb{C}_+)$, that is,
$\|\mathcal{L}f\|_{H^2(\mathbb{C}_+)}\approx\|f\|_{L^2(0,\infty)}$ (see \cite{LOLS}, for instance).

\smallskip

Let $M$ denote the involution
$$M(s)=\frac{1-s}{1+s}, \qquad (s\in \CC_+)$$
between $\CC_+$ and $\DD$. We have that $f\in H^2(\DD)$ if and only if $Vf \in H^2(\CC_+)$ with an equivalence of norms, where
$$Vf(s)= \frac{1}{1+s} f(M(s)), \qquad (s \in \CC_+)$$
(see, e.g. \cite[Chap. 8]{hoffman}). We summarize the change of coordinates as follows.

\begin{prop}
Let $(z_k)_{k\geq 0}$ be a sequence in $\DD$ and $(c_k)_{k\geq 0}$ a sequence of complex numbers. Then given a sequence $(G_k)$ of linear mappings $G_k: \CC^N \to \CC$, $N\geq 1$, the function $f \in H^2(\DD)$ satisfies the interpolation conditions
$$G_k f(z_k)=c_k$$
for each $k$ if and only if the function $F=Vf$ in $H^2(\CC_+)$ satisfies the interpolation condition
$$G_k (Vf)(s_k)=c_k$$
where $s_k=M(z_k)$ for each $k$.
\end{prop}

Let $H^2(\mathbb{D}; \mathbb{C}^N)$ denote the space of $\mathbb{C}^N$-valued functions $f = (f_1,\ldots,f_N)$ that are analytic on $\mathbb{D}$
and $\|f\|_{H^2(\mathbb{D}; \mathbb{C}^N)}^2:= \sum_{k=1}^N \|f_k\|_{H^2(\mathbb{D})}^2$, or equivalently,
\[
\|f\|_{H^2(\mathbb{D}; \mathbb{C}^N)}^2= \sup_{0<r<1} \frac{1}{2\pi}
\int_{0}^{2\pi} \|f(re^{i\theta})\|_{\mathbb{C}^N}^2 \, d\theta
< \infty,
\]
where $\|\cdot\|_{\mathbb{C}^N}$ is the Euclidean norm on $\mathbb{C}^N$. As a consequence, the following conditions are equivalent:


\begin{itemize}
  \item $(G_k)_{k\geq 0}$ and $(z_k)_{k\geq 0}$ have the property that for all $(c_k)_{k\geq 0} \in \ell^2$
  we can find $f \in H^2(\DD; \CC^N)$ such that $$G_k f(z_k)=c_k$$ for all $k$;
  \item $(G_k)_{k\geq 0}$ and $(s_k)_{k\geq 0}=(M(z_k))_{k\geq 0}$ have the property that for all $(c_k)_{k\geq 0} \in \ell^2$
  we can find $F \in H^2(\CC_+; \CC^N)$ such that $$(1+s_k)G_k F(s_k)=c_k$$ for all $k$.
  \end{itemize}


Now, Theorem 2.24 of \cite{JPP07}, combined with Theorem 2 of \cite{JZ01}, gives the following result for scalar interpolation in $H^2(\CC_+;\CC^N)$.
We write $\EE(F)=(G_k F(s_k))$ for $F \in H^2(\CC_+;\CC^N)$, where each $G_k$
is a linear mapping from $\CC^N$ to $\CC$, and thus has the form $G_k(x)=\langle x, g_k \rangle$ for some $g_k \in \CC^N$.


\begin{thm}
 The following are equivalent:
 \begin{itemize}
 \item $\EE(H^2(\CC_+;\CC^N))= \ell^2$;
 \item $(s_k)$ is the union of at most $N$ Carleson sequences, the
 quantities\\ $\|g_k\|^2/\re s_k$ are bounded above and below, and there exists a constant
 $r>0$ such that
 \[
 \inf_{m \in \NN} \min_{s_n \in \Lambda_m(r)} \angle (g_n e^{-s_n t}, \spam_{s_j\in\Lambda_m(r), j \ne n}
 \{ g_j e^{-s_j t}\} ) > 0,
 \]
 where $\Lambda_m(r) = \{ s_n : \left| \dfrac{s_n-s_m}{s_n+\overline{s_m}} \right| < r \}$
 and the angle is calculated in $L^2(0,\infty; \CC^N)$.
 \item The set $\{ g_k e^{-s_k t}: k \in \NN\}$ is a Riesz basic sequence
 in $L^2(0,\infty; \CC^N)$.
 \end{itemize}
 \end{thm}

 Note that the angle may  be expressed in terms of norms of operators, as in \cite[Lem.~2.15]{JPP07}.

On translating this back to the disc, noting that
$$\frac{1}{1+s_k}=\frac{1+z_k}{2}\text{ and } \re s_k= (1-|z_k|^2)/|1+z_k|^2,$$
and having in mind that the norms in $H^2(\CC_+)$ and  $L^2(0,\infty)$  are equivalent by the Laplace
transform, together with the properties of $V$ above, we obtain the following,
on writing $ \widetilde \EE(f)=(G_k f(z_k))$ for $f \in H^2(\CC_N)$ and
$G_k(x)=\langle x, g_k\rangle$ again.

 \begin{thm}\label{thm:interpold}
 The following are equivalent:
 \begin{itemize}
 \item $\widetilde \EE(H^2(\CC_+;\DD^N))= \ell^2$;
 \item $(z_k)$ is the union of at most $N$ Carleson sequences, the
 quantities\\ $\|g_k\|^2 (1-|z_k|^2)$ are bounded above and below, and there exists a constant
 $r>0$ such that
 \[
 \inf_{m \in \NN} \min_{z_n \in \Lambda_m(r)} \angle (g_n(1+z)/(1-z_n z), \spam_{z_j\in\Lambda_m(r), j \ne n}
 \{ g_j(1+z)/(1-z_j z)\} ) > 0,
 \]
 where $\Lambda_m(r) = \{ z_n : \left| \dfrac{z_n-z_m}{z_n - \overline{z_m}} \right| < r \}$
 and the angle is calculated in $H^2(\DD; \CC^N)$.
 \item The set $\{g_n(1+z)/(1-z_n z): k \in \NN\}$ is a Riesz basic sequence
 in $H^2(\DD; \CC^N)$.
 \end{itemize}
 \end{thm}


For $N=1$, where the $g_k$ are scalars, this reduces to saying that $(z_k)$ is a Carleson sequence
and the $|g_k|/\sqrt{1-|z_k|}$ are bounded above and below (the condition on angle is vacuous if
the $(z_k)$ are a Carleson sequence since $|z_n - \overline{z_m}|<|1-\overline{z_m}z_n|$). This agrees with McPhail's
Theorem \ref{thm:mcphail}.
\\


\subsection*{An explicit expression: the bivariable case} Let us write out in detail the interpolation question linked with
the condition for $\{T^n u,T^n v: n \ge 0\}$ to be  a frame.

The mapping
\[
\begin{array}{rcl}
\ell^2(\mathbb{C}^2) & \longrightarrow & H \\[4pt]
(a_n,b_n)_n & \longmapsto & \displaystyle \sum_{n} a_n T^n u + b_n T^n v
\end{array}
\]
should be bounded and surjective. Writing $f(z)=\sum_{k=0}^\infty a_k z^k$ and $g(z)=\sum_{k=0}^\infty b_k z^k$, with $f,g \in H^2$, we
have
\[
\sum_{n=0}^\infty  a_n T^n u+b_n T^n v = \sum_{k=0}^\infty f(\lambda_k)u_k z^k + g(\lambda_k) v_k z^k,
\]
so the interpolation problem is
$$f(\lambda_k)u_k+  g(\lambda_k)v_k=c_k,$$
for every $k\geq 0$, which fits into the framework of Theorem~\ref{thm:interpold}.\\

\section{Carleson frames and the M\"untz--Sz\'asz condition}\label{sec:3}

As pointed out in the Introduction, Carleson frames (or Alcamota frames), namely, frames satisfying conditions (1)-(4) in Theorem \ref{thm:diagonal}, exhibit excessive redundancy. Motivated by the recent work  in \cite{ChristensenHasannasabPhilippStoeva2024} and \cite{KrishtalMiller2026}, the authors in \cite{ACKM25} conjectured that if $b\in\ell^{2}$, $\sigma(D)\subset (0,1)$ the orbit $\{D^{\ell}b:\ell=0,1,\dots\}$ forms a Carleson frame for $\ell^{2}$ and  $\Lambda \subset (0,\infty)$ is a set satisfying the M\"untz--Sz\'asz condition
\begin{equation}\label{MS condition-2}
\sum_{\lambda \in \Lambda} \frac{1}{\lambda} = \infty,
\end{equation}
then the system $\{D^{\lambda} b : \lambda \in \Lambda\}$ is a frame for $\ell^2$.}

\begin{rem}{\rm
It is easy to see that \eqref{MS condition-2} by itself is not sufficient. Indeed, let $D$ be a diagonal operator such that $D\varphi_k=\mu_k \varphi_k$ for $\{\varphi_k\}_{k\geq 0}$ an
orthonormal basis of an infinite dimensional separable complex Hilbert space $H$. Suppose that $\{D^n u:\; n\geq 0\}$ is a frame for some $u \in H$, so that $\{\mu_k\}_{k\geq 0}$ is a Carleson sequence; we shall also assume that $0< \mu_k < 1$ for each $k$.


Now take a sequence $\Lambda=\{\lambda_n\}_{n\geq 0}$ such that $ \lambda_n>0$ for each $n$ and
$\lambda_n \to 1$. We will take the $\lambda_n$ distinct, and we see that
$\sum_{n=0}^\infty 1/\lambda_n = \infty.$ (Later we shall see the stronger result that even taking {\em subsequences of $\NN$}
satisfying the M\"untz--Sz\'asz condition is not enough to guarantee a frame.)


We claim that for every vector $v \in H$, say, $v = \sum_{k=0}^\infty v_k \varphi_k$, the
sequence $\{D^{\lambda_n} v:\; n\geq 0\}$ is not a frame. Clearly we may assume that $v \ne 0$, so let $k$ be chosen such that $v_k \ne 0$.
Now $\langle \varphi_k, D^n v \rangle_H = \ol {v_k} \mu_k^{\lambda_n}$ and
so
\[
\sum_{n=0}^\infty |\langle \varphi_k, D^n v \rangle_H| =\sum_{n=0}^\infty |v_k| \mu_k^{\lambda_n} = \infty,
\]
since $\mu_k^{\lambda_n} \to \mu_k$.


}
\end{rem}



Our next result states establishes that $\{D^{\lambda}b:\lambda\in\Lambda\}$ is a frame for $\ell^2$ if and only if $\{t^{\lambda}\}_{\lambda\in\Lambda}$ is a frame in a associated $L^2(\nu)$ space. Recall that if  \(X\) is a measurable space and \(x \in X\),  the \emph{delta measure} (or \emph{Dirac measure}) at \(x\), denoted by \(\delta_x\), is the measure defined by
\[
\delta_x(A) =
\begin{cases}
1, & \text{if } x \in A, \\
0, & \text{if } x \notin A,
\end{cases}
\]
for every measurable set \(A \subseteq X\). Equivalently, $\delta_x(A) = \mathbf{1}_A(x),$ where \(\mathbf{1}_A\) is the indicator function of the set \(A\).


\begin{thm}\label{prop-ms}
Let $D$ be a diagonal operator on $\ell^2$ with spectrum
$\sigma(D)=\{\mu_k:k\ge 0\}\subset(0,1)$ and $\{\phi_k\}_{k\geq 0}$ an orthonormal basis such that $D\phi_k=\mu_k\phi_k$ for every $k$.
Let $b\in\ell^2$ and assume that $\{D^{\ell} b:\ell=0,1,2,\dots\}$ is a Carleson frame for $\ell^2$. If $\Lambda=\{\lambda_k\}_{k\geq 0}$ is any sequence of non-negative real numbers, then $\{D^\lambda b:\lambda\in\Lambda\}$
is a frame for $\ell^2$ if and only if $\{t^{\lambda}\}_{\lambda\in\Lambda}$ is a frame in $L^2(\nu)$ where $\nu$ is the discrete measure
\[
\nu= \sum_{k\ge 0} (1-\mu_k^2)\,\delta_{\mu_k}
\quad\text{on } [0,1).
\]
\end{thm}




The proof is based on an spectral model in which $D$ is represented as multiplication by $t\in(0,1)$ on a Hilbert space determined by $(D,b)$.
In this model, the vectors $D^{\lambda}b$ correspond precisely to the functions $t^{\lambda}$.


\begin{proof}
Let us write $b=\sum_k b_k\phi_k$ and note that, since $\{D^\ell b\}_{\ell\ge 0}$ is a Carleson frame and $\sigma(D)\subset(0,1)$,
the frame characterization for diagonal normal operators yields:
\begin{itemize}
\item $(\mu_k)$ is an interpolating sequence in $\mathbb D$,
\item there exist constants $C_1,C_2>0$ such that
\begin{equation}\label{uk}
C_1 \le \frac{|b_k|}{\sqrt{1-\mu_k^2}} \le C_2
\quad\text{for all } k.
\end{equation}
\end{itemize}
For $x=\sum_{k\geq 0} x_k\phi_k\in\ell^2$ and $\lambda>0$, note that
\[
\ip{x}{D^\lambda b}
= \sum_{k\ge 0} x_k\overline{b_k}\,\mu_k^\lambda.
\]
Define $J:\ell^2\to L^2(\nu)$ by
\[
(Jx)(\mu_k) := \frac{x_k\overline{b_k}}{1-\mu_k^2}.
\]
Using \eqref{uk}, $J$ is bounded and bounded below, and
\[
\norm{Jx}_{L^2(\nu)}^2
= \sum_{k\geq 0} \frac{|b_k|^2}{1-\mu_k^2}|x_k|^2
\asymp \norm{x}_{\ell^2}^2.
\]
Moreover,
\[
\ip{x}{D^\lambda b}
= \sum_{k\geq 0} (1-\mu_k^2)(Jx)(\mu_k)\,\mu_k^\lambda
= \ip{Jx}{t^\lambda}_{L^2(\nu)}.
\]
Hence, $\{D^\lambda b\}_{\lambda\in\Lambda}$ is a frame for $\ell^2$ if and only if $\{t^\lambda\}_{\lambda\in\Lambda}$ is a frame for $L^2(\nu)$, as we claimed.
\end{proof}


With Theorem \ref{prop-ms} at hand, in order to approach the conjecture in \cite{ACKM25}, our task is reduced to studying when $\{t^\lambda\}_{\lambda\in\Lambda}$ is a frame for $L^2(\nu)$ assuming that $\{t^\ell\}_{\ell\in\mathbb{N}}$ is. For such a goal, let us write for the moment $w_k = 1 - \mu_k^2$ for $k\geq 0$. Note that the frame inequality would require
\[
A \|f\|_{L^2(\nu)}^2
\le
\sum_{\lambda \in \Lambda}
|\langle f, t^\lambda \rangle|^2.
\]
Testing on the normalized atom
\begin{equation}\label{ek}
e_k =
\frac{\mathbf{1}_{\{\mu_k\}}}{\sqrt{w_k}},
\end{equation}
we get
$
\|e_k\|_{L^2(\nu)} = 1
$
and
\[
\langle e_k, t^\lambda \rangle
=
\sqrt{w_k}\,\mu_k^\lambda.
\]
Therefore a necessary condition for the lower frame bound is
\[
\inf_{k\geq 0}
w_k
\sum_{\lambda \in \Lambda}
\mu_k^{2\lambda}
> 0.
\]
That is,
\[
\inf_{k\geq 0}
(1 - \mu_k^2)
\sum_{\lambda \in \Lambda}
\mu_k^{2\lambda}
> 0.
\]
In general, $\{t^\lambda\}_{\lambda \in \Lambda}$ is a frame for $L^2(\nu)$ exactly when there are constants $A,B > 0$ such that, for every finitely supported scalar sequence $(a_k)$,
\[
A \sum_{k\geq 0} |a_k|^2 w_k
\le
\sum_{\lambda \in \Lambda}
\left|
\sum_{k\geq 0} a_k w_k \mu_k^\lambda
\right|^2
\le
B \sum_{k\geq 0} |a_k|^2 w_k.
\]
After normalizing \(c_k = a_k \sqrt{w_k}\),
\[
A \sum_{k\geq 0} |c_k|^2
\le
\sum_{\lambda \in \Lambda}
\left|
\sum_{k\geq 0} c_k \sqrt{1 - \mu_k^2}\, \mu_k^\lambda
\right|^2
\le
B \sum_{k\geq 0} |c_k|^2.
\]
So, a necessary pointwise consequence is
\begin{equation}\label{pointwise condition}
0 < \inf_{k\geq 0} (1 - \mu_k^2)\sum_{\lambda \in \Lambda} \mu_k^{2\lambda}
\le
\sup_{k\geq 0} (1 - \mu_k^2)\sum_{\lambda \in \Lambda} \mu_k^{2\lambda}
< \infty.
\end{equation}
Indeed, \eqref{pointwise condition} allows us to show that even when $\Lambda=\{\lambda_n\}_{n\geq 0}\subset \mathbb{N}$,  \eqref{MS condition-2} does not suffice to prove that  $\{D^{\lambda}b:\lambda\in\Lambda\}$ is a frame for $\ell^{2}$.
First, we need an estimate.

\begin{lem}\label{lem:xsx}
For $x>0$ let
\[
S(x)= \sum_{n=2}^\infty \exp(-nx\log(n)) .
\]
Then  $\lim_{x \to 0+} xS(x) =0$.
\end{lem}
\begin{proof}
Fix $x \in (0,1/4)$ and  take $N=\lfloor 1/(x \log (1/x)) \rfloor$.
Then
\[
S(x)= \sum_{n=2}^N \exp(-nx\log(n)) + \sum_{n=N+1}^\infty \exp(-nx\log(n))= S_1(x)+S_2(x), \quad \hbox{say}.
\]
For $2\le n\le N$, we have $\exp(-nx\log(n))\le 1$, so
$S_1 (x)\le N\le 1/(x \log (1/x))$. For the tail of the series,
\[
S_2(x) \le \int_{N-1}^\infty\exp(-xt\log t) \, dt.
\]
We take $u=xt \log t$ so that
\[
du=x(\log t+1)\,  dt \ge x \log t \, dt \ge x \log  (N-1) \, dt
\]
and we have as an upper bound
\[
S_2(x) \le \frac{1}{x \log (N - 1)}\int_{0}^\infty \exp(-u) \ du,
\]
so that $S_2(x) \le 1/(x \log (N-1)) \le C/(x \log (1/x))$ for some $C>0$ independent of $x$.

Putting these together, we see that $xS(x)$ is bounded by a constant times $1/\log(1/x)$ as $x \to 0+$, and hence
it tends to $0$.
\end{proof}

\begin{exm}
If we now take $\lambda_n=\lceil n \log n \rceil$ or $\lambda_n=p_n$, the $n$th prime, for $n \ge 2$, both of which are greater than
$n \log n$ (in the case of $p_n$ this is Rosser's theorem~\cite{rosser}), then
each sequence satisfies the M\"untz--Szász  condition $\sum_{n=2}^\infty 1/\lambda_n = \infty$ and, by Lemma~\ref{lem:xsx}, the condition
\[
\lim_{x \to 0+} x\sum_{n=2}^\infty \exp(-\lambda_n x)=0 .
\]
Now if $\mu_k \to 1-$, then
\[
\inf_k
(1-\mu_k^2)
\sum_{n=2}^\infty
\mu_k^{2\lambda_n}
=0,
\]
since, with $x=1-\mu_k^2$ small, we have
$\mu_k^2 = 1-x \le \exp(-x)$.


Then testing the frame inequality on $e_k$ (defined in \eqref{ek}) gives
\[
\sum_{n\geq 2} |\langle e_k, t^{\lambda_n} \rangle|^2
=
(1-\mu_k^2)\sum_{n\geq 2} \mu_k^{2\lambda_n}
\to 0,
\]
while
\[
\|e_k\|_{L^2(\nu)} = 1.
\]
So no positive lower frame bound can hold.
\end{exm}

As a byproduct of our previous discussion, we provide a necessary condition for $\{t^\lambda\}_{\lambda \in \Lambda}$ to be a frame for $L^2(\nu)$ whenever $\{t^\ell\}_{\ell \in \mathbb{N}}$ is.

\begin{proposition}
Let $\{\mu_k\}_{k\geq 0}\subset(0,1)$ be an interpolating sequence for $H^2(\mathbb{D})$, and let $\nu$ the measure on $[0,1)$
\[
\nu=\sum_{k\ge0}(1-\mu_k^2)\delta_{\mu_k}.
\]
Assume that $\{t^\ell\}_{\ell \in \mathbb{N}}$ is a frame in $L^2(\nu)$ and $\Lambda=\{\lambda_n\}_{n\ge0}\subset (0,+\infty)$. If the family $\{t^{\lambda_n}\}_{n\geq 0}$ is a frame for $L^2(\nu)$ then both conditions:
\begin{enumerate}
\item \textbf{the Müntz--Szász condition}
\[
\sum_{n\ge0}\frac1{\lambda_n}=\infty,
\]
and
\item \textbf{a Carleson-type condition}
\[0<
\inf_{k\ge0}
(1-\mu_k^2)\sum_{n\ge0}\mu_k^{2\lambda_n}
\le
\sup_{k\ge0}
(1-\mu_k^2)\sum_{n\ge0}\mu_k^{2\lambda_n}
<\infty,
\]
\end{enumerate}
hold.
\end{proposition}


\begin{proof}
Let us assume that $\{t^{\lambda_n}\}_{n\geq 0}$ is a frame for $L^2(\nu)$. Then it is complete, and therefore the closed linear span of
$\{t^{\lambda_n}:n\geq 0\}$ is all of $L^2(\nu)$. Arguing by contradiction, if the Müntz--Szász condition does not hold, namely,
$\sum_{n\geq 0}\frac1\lambda_n<\infty,$ then the Müntz--Szász theorem implies that
$$
\overline{\operatorname{span}}\{t^{\lambda_n}:n\geq 0\}
$$
is not dense on any infinite subset of $[0,1)$ with an accumulation point in $[0,1]$ away from $0$. Since the support of $\nu$ is infinite and has accumulation point $1$, this contradicts completeness in $L^2(\nu)$ and therefore, the Müntz--Szász condition is necessary.


Next, let \(A,B>0\) be frame bounds for \(\{t^{\lambda_n}\}_{n\ge0}\) in $L^2(\nu)$. Arguing as in \eqref{pointwise condition} for
$e_k=\frac1{\sqrt{w_k}}\mathbf 1_{\{\mu_k\}}$, $k\geq 0$, it follows
\[
0<
\inf_{k\ge0}
w_k\sum_{n\ge0}\mu_k^{2\lambda_n}
\le
\sup_{k\ge0}
w_k\sum_{n\ge0}\mu_k^{2\lambda_n}
<\infty,
\]
which yields the statement.
\end{proof}



\begin{remark}
Let us denote \(w_k=1-\mu_k^2\) for \(k\ge0\), so
$\nu=\sum_{k\ge0}w_k\delta_{\mu_k},
$
and the space \(L^2(\nu)\) is naturally identified with \(\ell^2\) by the unitary map
$$\begin{array}{lccc}\label{def U}
U:&L^2(\nu)&\to&\ell^2,\\
& f& \to &  Uf=\{f_k\}_{k\ge0}
\end{array}
$$
where
\begin{equation}\label{def U}
f_k=\sqrt{w_k}\,f(\mu_k), \quad (k\ge0).
\end{equation}
Indeed,
\[
\|Uf\|_{\ell^2}^2
=
\sum_{k\ge0}|f_k|^2
=
\sum_{k\ge0}|f(\mu_k)|^2w_k
=
\|f\|_{L^2(\nu)}^2.
\]

\noindent For \(\lambda_n\in\Lambda\), we have
\[
\langle f,t^{\lambda_n}\rangle_{L^2(\nu)}
=
\sum_{k\ge0}f(\mu_k)\mu_k^{\lambda_n}w_k
=
\sum_{k\ge0}f_k \sqrt{w_k} \mu_k^{\lambda_n}.
\]
Hence the frame problem in \(L^2(\nu)\) for \(\{t^{\lambda_n}\}_{n\ge0}\) is equivalent to the frame problem in \(\ell^2\) for the vectors
\begin{equation}\label{frame problem}
b_{\lambda_n}
=
\left(\sqrt{w_k}\mu_k^{\lambda_n}\right)_{k\ge0}, \qquad  (n\geq 0).
\end{equation}
Indeed,
\begin{equation}\label{frame problem--}
\langle f,t^{\lambda_n}\rangle_{L^2(\nu)}
=
\langle Uf,b_{\lambda_n}\rangle_{\ell^2}.
\end{equation}
\end{remark}



\section{Iterated multiplication and weighted composition operators}\label{sec:4}


Take $u \in H^2(\DD)$ and $\phi \in H^\infty$.
Cheshmavar~\cite{chesh} posed the question when the sequence $(\phi^n u)_{n\geq 0}$
forms a frame in $H^2(\DD)$, but did not give a complete answer. We shall solve this problem using weighted composition operators.

It was assumed that $\|\phi\|_\infty \le 1$, and indeed this follows from the uniform boundedness theorem,
since if we have a frame then $(\langle \phi^n u, v \rangle)$ is a bounded sequence for every
$v \in H^2$ and hence $(\phi^n u)$ is bounded in norm.

Now we use the condition given in \eqref{eq:bddsurj}, and this says that the
mapping $T^*: \ell^2 \to H$ should be bounded and surjective. Writing this
as a  mapping     $H^2 \to H^2$ it says that the operator given by
\[
W: \sum_{n=0}^\infty a_n z^n \mapsto \sum_{n=0}^\infty a_n \phi^n u
\]
should be bounded and surjective. But $W$ is simply the weighted composition operator $W_{\phi,u}:=M_u C_\phi$.

The following result is given by Gunatillake \cite[Thm.~2.0.1]{gun}, written here in the
notation used above. Note that $W_{\phi,u}$ is injective (except in   trivial cases)
and so boundedness and surjectivity together are equivalent to invertibility.

\begin{thm}
The operator $W_{\varphi, u}$ on $H^2(\DD)$, defined by
\[
   (W_{\varphi, u} f)(z) =  u(z) f(\varphi(z)),
\]
 is invertible on $H^2(\DD)$ if and only if $\phi$ is an automorphism of $\DD$
and $u$ is both bounded and bounded away from zero on $\DD$.
\end{thm}





Write $\mathbf{K}_w(z)=\frac{1}{1-\overline{w}z}$ for the reproducing kernel in $H^2(\DD)$, so
$f(w)=\langle f,\mathbf{K}_w\rangle$.
It is well known that for
 every $w\in\mathbb D$,
\[
W_{\varphi, u}^*\,\mathbf{K}_w \;=\; \overline{ u(w)}\,\mathbf{K}_{\varphi(w)}.
\]
With Parseval frames in mind, let us recall the following theorem by Bourdon and Narayan \cite{BourdonNarayan2010}.

\begin{thm}\label{thm:BN}
The operator $W_{\varphi, u}$ on $H^2(\DD)$, defined by
\[
   (W_{\varphi, u} f)(z) =  u(z) f(\varphi(z)),
\]
is unitary if and only if $\varphi$ is a disk automorphism and
\[
    u(z) = c \, \frac{\mathbf{K}_p(z)}{\|\mathbf{K}_p\|},
\]
where $\varphi(p) = 0$, $\mathbf{K}_p$ is the reproducing kernel at $p$, and $|c| = 1$.
\end{thm}





\begin{prop}[Explicit evaluation formula]\label{prop:pointwise-adjoint}
For $f\in H^2$ and $w\in\mathbb D$,
\[
(W_{\varphi, u}^* f)(w)
=\big\langle f,\, u(\cdot)\,\mathbf{K}_w(\varphi(\cdot))\big\rangle_{H^2}
=\int_{0}^{2\pi}
f(e^{it})\,\overline{ u(e^{it})}\,
\frac{1}{1-w\,\overline{\varphi(e^{it})}}\;\frac{dt}{2\pi}.
\]
\end{prop}

\begin{proof}
By the reproducing property,
\[
(W_{\varphi, u}^* f)(w)
=\langle W_{\varphi, u}^* f,\,\mathbf{K}_w\rangle
=\langle f,\,W_{\varphi, u}\mathbf{K}_w\rangle
=\langle f,\, u(\cdot)\,\mathbf{K}_w(\varphi(\cdot))\rangle,
\]
and the boundary integral representation of the $H^2(\DD)$ inner product yields the second equality.
\end{proof}

Cowen gave the following adjoint formula for linear-fractional symbols (see \cite[Thm.~9.2]{cowenmaccluer}).

\begin{thm} \label{thm:cowen}
Suppose $\varphi$ is linear-fractional,
\[
\varphi(z)=\frac{az+b}{cz+d},\qquad ad-bc\neq 0,
\]
and define
\[
\sigma(z)=\frac{\overline{a}\,z-\overline{c}}{-\overline{b}\,z+\overline{d}},
\qquad
h(z)=cz+d,
\qquad
g(z)=\frac{1}{-\overline{b}\,z+\overline{d}}.
\]
Then for the composition operator $C_\varphi$ on $H^2(\DD)$,
\[
C_\varphi^* \;=\; M_g\,C_\sigma\,M_h^*.
\]
\end{thm}


\begin{cor}[Unitary case]\label{cor:unitary}
Assume $\varphi$ is a disk automorphism and $ u(z)=c\,\dfrac{\mathbf{K}_p(z)}{\|\mathbf{K}_p\|}$ with $\varphi(p)=0$ and $|c|=1$. Then $W_{\varphi, u}$ is unitary, and
\[
W_{\varphi, u}^* = W_{\varphi, u}^{-1} = M_{\eta}\,C_{\varphi^{-1}},
\qquad
\eta(z)=\frac{1}{ u(\varphi^{-1}(z))}\quad (\text{unimodular a.e.}).
\]
\end{cor}

\begin{proof}
For a disk automorphism $\varphi$, $C_\varphi$ is unitary on $H^2(\DD)$; the chosen $ u$ is unimodular a.e., so $M_ u$ is unitary. Hence $W_{\varphi, u}$ is unitary and its adjoint equals its inverse. The explicit inverse follows from
\[
W_{\varphi, u}\,M_{\eta}\,C_{\varphi^{-1}}
= M_ u\,C_\varphi\,M_{\eta}\,C_{\varphi^{-1}}
= M_{ u\,(\eta\circ \varphi)}\,I,
\]
and choosing $\eta=1/( u\circ\varphi^{-1})$ gives the identity operator.
\end{proof}

It is of independent interest to consider the question when the adjoint of a weighted composition operator is an isometry (not necessarily unitary). One obvious necessary condition is
that it is isometric on reproducing kernels $\mathbf{K}_w$ for $w \in \DD$, where
\[
\mathbf{K}_w(z)= \frac{1}{1-\overline w z} \qquad (z \in \DD).
\]
It is well known and easily verified that
\[
W_{\phi, u}^* \mathbf{K}_w = \overline{ u(w)} \mathbf{K}_{\phi(w)}
\]
and so a necessary condition is that
\begin{equation}\label{eq:psiw}
| u(w)|^2 = \frac{1-|\phi(w)|^2}{1-|w|^2} \qquad (w \in \DD).
\end{equation}
Since operators with isometric adjoints are surjective (so in this case, invertible) we
must also have $\phi$ an automorphism and $ u$ both bounded and bounded away from zero.

Now if $\phi$ is an automorphism with a zero at $p$, then by an easy calculation this condition \eqref{eq:psiw} becomes
\[
| u(w)|^2 = \frac{1-|p|^2}{|1-\overline p w|^2} \qquad (w \in \DD),
\]
which, since $ u$ is analytic, implies by Theorem~\ref{thm:BN} that $W$ and $W^*$ are
unitary.

This can also be expressed by saying that we have a reproducing kernel hypothesis for
the isometry condition of an adjoint of a weighted composition operator.

We therefore have the following result.

\begin{thm}
The adjoint of a weighted composition operator on $H^2(\DD)$ is an isometry if and only
if it is unitary.
\end{thm}




\section{Model spaces and their generalizations}\label{sec:6}

Let $\theta \in H^\infty$ be an inner function. The associated model space is
defined by
\[
K_\theta= H^2 \ominus \theta H^2,
\]
the orthogonal complement of the shift-invariant subspace $\theta H^2$ in the
Hardy space $H^2$. By Beurling's theorem \cite{Beu}, every nontrivial closed subspace of $H^2$ invariant under
the unilateral shift has the form $\theta H^2$ for some inner function $\theta$.
Consequently, the spaces $K_\theta$ are invariant under the backward shift operator
\[
S^*f(z)=\frac{f(z)-f(0)}{z},
\]
and the nontrivial invariant subspaces of $S^*$ are precisely the model spaces
$K_\theta$. A good reference for model spaces is the book~\cite{GMR}.

For $w \in \DD$ we shall write $k^\theta_w \in K_\theta$ for the reproducing kernel
for evaluation at $w$. That is $\langle f, k^\theta_w \rangle = f(w)$ for all $f \in K_\theta$.

\subsection{The case of a single operator}



The paper \cite{CHP20}
contains the following observations. Note that saying $(A_1,u_1)$ is similar to $(A_2,u_2)$ means that there is an invertible $V$
such that
$A_2=V A_1 V^{-1}$ and $u_2=A u_1$.

\begin{itemize}
\item The system $(A^n u)_{n =0}^\infty $ forms a frame if and only if $(A,u)$ is similar to
either\\ (i) the compressed shift
$(S_\theta, k_0^\theta)$ acting on the model space $K_\theta$,
with $k_0^\theta$ the reproducing kernel at $0$, that is, $k_0^\theta(z)=1-\ol{\theta(0)}\theta(z)$,
or\\ (ii) the standard shift   $(S, e_0)$ acting on $H^2(\DD)$, with $e_0(z)=1$.
\item In the first case, $S_\theta$ is a completely non-unitary $C_0$ contraction, and
$\theta$ is its minimal function (unique to within a unimodular constant). That is,
$\theta$ is an inner function and
$\theta(S_\theta)=0$.
\end{itemize}

Note that we have an identity for the spectrum $\sigma(A)=\sigma(S_\theta)=\sigma(\theta)$
(the Liv\v{s}ic--Moeller theorem \cite{nikolski-shift}) in case (i), and $\sigma(A)=\ol\DD$ in case (ii), although this does not
determine $\theta$ uniquely in case (i), unlike the minimal function.

\subsection{Particular cases}\label{subsec:pc}

\begin{enumerate}
\item For the frame $(A,u)$ with $A$ a diagonal operator,   discussed in Section~\ref{sec:2}, we either have $A$ similar to
a compressed shift $S_\theta$ or to
the shift $S$.
In the first case the operator is $C_0$ and its eigenvalues must form a Blaschke sequence
(clearly with no repeats). Further,
if the normalized eigenvectors are to form a Riesz basis then the eigenvalues are
a Carleson sequence.
The second case is impossible since $S$ has no eigenvalues.

\item The case when $A$ is similar to a direct sum of Jordan blocks
where the generalized eigenvectors form a Riesz basis
 corresponds to an inner function $\theta$ with repeated roots.

Observations:
\begin{itemize}
\item We can only have one block for each  eigenvalue, since $f \in \ker(S_\theta-\lambda I)$
gives $zf=\theta h + \lambda f$ for some $h \in H^2$ and $f=\theta h/(z-\lambda)$.
If this is in $K_\theta$ we have $\theta(\lambda)=0$ and, dividing by $\theta$, $h/(z-\lambda) \in \ol{H^2_0}$ so $h$ is a constant. That is, the eigenspace is one-dimensional, being spanned by $\theta/(z-\lambda)$.

It can be seen by considering the Taylor series
that the eigen\-spaces of $S_\theta^*$ are also one-dimensional, the eigenvectors being reproducing kernels, and the eigenvalues are $\ol\lambda$, where $\theta(\lambda)=0$.

\item Also, $S_\theta$ is a $C_0$ operator, with minimal function $\theta$, which, since
each zero of
$\theta$ has finite order, means that the blocks are all finite.

\item It does not seem to be known whether generalized eigenvectors can form a
Riesz sequence (if the zeros form a Carleson set and only finitely-many zeros are repeated,
then this will happen), so we do not know whether it is possible to have infinitely-many
nontrivial blocks if there is a frame.

\end{itemize}

\item
For the multiplication frame $(M_\phi, f)$ discussed in Section~\ref{sec:3},
we see that with $\phi$ an automorphism we are in situation (ii) above (similarity with the shift); the more complicated situation (i) does not occur.\\

\item As discussed in \cite{GPR-Hardy}, on writing $\mathcal{H}_1$ for the Hardy operator
on $L^2(0,1)$ defined by
$$(\mathcal{H}_1 f)(x) = \frac{1}{x} \int_{0}^{x} f(t) dt, \quad 0 < x < 1,$$
the fact that $I-\mathcal{H}_1^*$ is unitarily equivalent to the
unilateral shift $S$ on $H^2(\DD)$ implies that it can generate frames. This is situation (ii) once more.
\end{enumerate}

\subsection{Parseval frames}




Starting with $(S, e_0)$ acting on $H^2(\DD)$ we have $\langle f, S^n e_0 \rangle = \hat f(n)$,
the $n$th Taylor coefficient,
so this is clearly a Parseval frame.

Moreover, for  $(S_\theta, k_0^\theta)$ acting on $K_\theta$ we have
\[
\langle f, S_\theta^n k_0^\theta \rangle = \langle (S^*)^n f, k_0^\theta\rangle = \hat f(n)
\]
once more, since $k_0^\theta$ is the reproducing kernel at $0$  for $K_\theta$.
Thus we have
\begin{cor}
Every frame of the form $(A^n u)_n$ is similar to a Parseval frame.
\end{cor}

A generalization of the model space approach described above is due to
Cabrelli, Molter and Su\'arez~\cite{CMS23}.
They give an answer to the following question:

{\em
Let $A$ be a bounded operator on a Hilbert space $H$ and $\{u_1,\ldots,u_m\}$ an
independent  set in  $H$. When is
$\{A^n u_k: n \ge 0, k=1,\ldots,m\}$   a frame?}

The answer can be expressed in terms of invariant subspaces of the iterated shift $S^m$
on $H^2$, and these in turn can be linked to the invariant subspaces of
the shift on the vector-valued Hardy space $H^2(\DD,\CC^m)$ by using the
 following commutative diagram:

\begin{align*}\begin{CD}
H^2 @> S^m  >>H^2\\
@V J   VV @ VV J \;V\\
H^2(\DD, \CC^m) @> S >> H^2(\DD, \CC^m)
\end{CD},  \end{align*}
where $J: H^2 \to H^2(\DD, \CC^m)$ is the unitary mapping
\[
J: \sum_{k=1}^m z^{k-1}f_k(z^m) \mapsto (f_1(z),f_2(z),\ldots, f_m(z)).
\]
for $f_1,\ldots, f_m \in H^2$ together with the
Beurling--Lax theorem (e.g. \cite[Thm.~3.1.7]{LOLS}).
Note that a function $\Phi \in H^\infty(\DD, \LL(\CC^r,\CC^m))$ is said to be {\em inner\/} if
$\Phi(e^{i\omega})$ is an isometry for almost all real $\omega$.


\begin{thm}\label{thm:kinvsub}
Let $K$ be a non-zero subspace of $H^2(\DD, \CC^m)$ that is invariant under the
shift $S$. Then there is an $r$ with $0 \le r \le m$ and an inner function
$\Phi \in H^\infty(\DD,  \LL(\CC^r,\CC^m))$ such that $K=\Phi H^2(\DD, \CC^r)$.
\end{thm}

That is, the subspaces  of $H^2(\DD)$ invariant under $S^m$ have the form $J^{-1}K$, where $K$ is
given in Theorem~\ref{thm:kinvsub}.
We next have that $\rang U^* = (\ker U)^\perp =: M$, a space invariant under $(S^*)^m$ and that $U:M \to H$
is an isomorphism.
Arguments similar to those in \cite{CHP20} now show that $A$  is similar to the compressed shift
$P_M(S^m_{|M})$. The simplest case is $M=H^2$, which makes $A$ similar to $S^m$.

Under the transformations given, we wish to find vectors corresponding to $u_1,\ldots,u_m$.
In the case $K=H^2(\DD,\CC^m)$
these will be $v_j= JU^{-1}u_j= e_{j}$, for $j=1,\ldots,m$, the constant function
that takes the value $(0,\ldots,0,1,0, \ldots 0)$ with the $1$ in the $j$th place.
In the more general case $K=\Phi H^2(\DD,\CC^r)$
it will be the projection of $v_j$ onto $K^\perp$, namely the vectorial reproducing kernel $k^M_{j,0}$
such that $\langle (f_1,\ldots,f_m),k^M_{j,0} \rangle = f_j(0)$ for
$(f_1,\ldots,f_m) \in M$. Once more these are Parseval frames.






\begin{thebibliography}{99}



\bibitem{ACNP25}
A. Aguilera, C. Cabrelli, F. Nagreira and V. Paternostro,
Optimal Dynamical Frames, Journal D'Analyse Mathematique, (2026) to appear.
{\tt https://arxiv.org/pdf/2506.00567}

\bibitem{ACCMP17}
A.~Aldroubi, C.~Cabrelli, A.F.~{\c{C}}akmak, U.~Molter, and A.~Petrosyan,
Iterative actions of normal operators,
\emph{J. Funct. Anal.} \textbf{272} (2017), 1121--1146.


\bibitem{ACMT17}
A. Aldroubi, C. Cabrelli, U. Molter and S. Tang, Dynamical sampling.
{\em Appl. Comput. Harmon. Anal.} 42 (2017), no. 3, 378--401.


\bibitem{ACKM25}
A.~Aldroubi, C.~Cabrelli, I.~Krishtal, and U.~Molter,
Dynamical Sampling: A Survey, La Matematica 5, 37 (2026).
%{\tt https://arxiv.org/abs/2511.10769}


\bibitem{Beu}
A.~Beurling,
\emph{On two problems concerning linear transformations in Hilbert space},
Acta Math. \textbf{81} (1949), 239--255.


\bibitem{BourdonNarayan2010}
P.S.~Bourdon and S.K.~Narayan,
Normal weighted composition operators on the Hardy space,
{\em J. Math. Anal. Appl.}  367  (2010), no.~1, 278--286.

\bibitem{CMPP20}
C. Cabrelli, U. Molter, V. Paternostro and F. Philipp, Dynamical sampling on finite index sets.
{\em J. Anal. Math.} 140 (2020), no. 2, 637--667.

\bibitem{CMS23}
C. Cabrelli, U. Molter and D. Su\'arez,   Frames of iterations and vector-valued model spaces.
{\em Sampling, approximation, and signal analysis--harmonic analysis in the spirit of J. Rowland Higgins}, 329--348, Appl. Numer. Harmon. Anal., Birkh{\"a}user/Springer, Cham, 2023.

\bibitem{Carleson}
L.~Carleson,
An interpolation problem for analytic functions,
{\em Amer.\ J.\ Math.}\ \textbf{80} (1958), 921--930.


\bibitem{chesh}
J. Cheshmavar, Frame properties induced by iteration of a multiplication operator on Hardy spaces.
{\tt https://arxiv.org/abs/2509.04879}

\bibitem{CHP20}
O. Christensen, M. Hasannasab and F. Philipp,   Frame properties of operator orbits.
{\em Math. Nachr.} 293 (2020), no. 1, 52--66.

\bibitem{ChristensenHasannasabPhilippStoeva2024}
O.~Christensen, M.~Hasannasab, F.M.~Philipp, and D.~Stoeva,
\newblock The mystery of Carleson frames,
\newblock \emph{Appl.\ Comput.\ Harmon.\ Anal.}, 72 (2024), Paper No.\ 101659, 5~pp.





\bibitem{cowenmaccluer}
C.C. Cowen and B.D. MacCluer,   {\em Composition operators on spaces of analytic functions}. Studies in Advanced Mathematics. CRC Press, Boca Raton, FL, 1995.



\bibitem{GPR-Hardy}
E.A. Gallardo-Guti\'errez, J,R. Partington and W,T. Ross,
Hardy operators: In the footsteps of Brown, Halmos, and Shields.
{\tt https://arxiv.org/abs/2501.14706}

\bibitem{GMR}
S.R. Garcia, J. Mashreghi and W.T. Ross, {\em Introduction to model spaces and their operators.} Cambridge Studies in Advanced Mathematics, 148. Cambridge University Press, Cambridge, 2016.

\bibitem{Garnett}
J.B.~Garnett,
\emph{Bounded Analytic Functions},
Revised first edition,
Graduate Texts in Mathematics, Vol.~236,
Springer-Verlag, New York, 2007.



\bibitem{gun}
G. Gunatillake,   Invertible weighted composition operators.
{\em J. Funct. Anal.} 261 (2011), no. 3, 831--860.

\bibitem{Hayman}
W.~Hayman,
Interpolation by bounded functions,
{\em Ann.\ Inst.\ Fourier (Grenoble)} \textbf{8} (1958), 277--290.



\bibitem{hoffman}
K. Hoffman,
{\em Banach spaces of analytic functions}.
Dover Publications, Inc., New York, 1988.

\bibitem{JPP07}
B. Jacob, J.R. Partington and S. Pott,   Interpolation by vector-valued analytic functions, with applications to controllability. {\em J. Funct. Anal.} 252 (2007), no. 2, 517--549.

\bibitem{JZ01}
B. Jacob and H. Zwart,   Exact observability of diagonal systems with a finite-dimensional output operator. {\em Systems Control Lett.} 43 (2001), no. 2, 101--109.



\bibitem{KrishtalMiller2026}
I.~Krishtal and B.~Miller,
\newblock Demystifying Carleson frames,
\newblock \emph{Applied and Computational Harmonic Analysis}, 80 (2026), Paper No.\ 101811, 8~pp.



\bibitem{mcphail}
J.D. McPhail,   A weighted interpolation problem for analytic functions.
{\em Studia Math.} 96 (1990), no. 2, 105--116.

\bibitem{meyer}
Y. Meyer, {\em Wavelets and operators}. Translated from the 1990 French original by D.L.~Salinger. Cambridge Studies in Advanced Mathematics, 37. Cambridge University Press, Cambridge, 1992.

\bibitem{Newman}
D.J.~Newman,
Interpolation in $H^\infty$,
{\em Transactions Amer. Math. Soc.} \ \textbf{92} (1958), 501--507.

\bibitem{nikolski-shift}
N.K. Nikol'ski\u \i,  {\em  Treatise on the shift operator. Spectral function theory.}  Grundlehren der mathematischen Wissenschaften, 273. Springer-Verlag, Berlin, 1986.


\bibitem{LOLS}
J.R. Partington, {\em Linear operators and linear systems.} London Mathematical Society Student Texts, 60. Cambridge University Press, Cambridge, 2004.


\bibitem{rosser} J.B. Rosser,  The $n$-th prime is greater than $n$\text{log}$n$,
{\em Proc.   London Math. Soc.} (2) 45 (1939), no. 1, 21--44.

\bibitem{ShapiroShields}
H.S.~Shapiro and A.L.~Shields,
On some interpolation problems for analytic functions,
{\em Amer.\ J.\ Math.}\ \textbf{83} (1961), 513--532.



\end{thebibliography}

\end{document}



